% ============================================================
%  GUÍA TAMARIN PROVER — ESTILO LIBRO
%  Tesis de Grado: Protocolo de Votación Electrónica PQC
%  Universidad Mayor de San Andrés (UMSA) — 2026
% ============================================================
\documentclass[12pt, a4paper, twoside, openright]{book}

% ── Codificación y lenguaje ─────────────────────────────────
\usepackage[utf8]{inputenc}
\usepackage[T1]{fontenc}
\usepackage[spanish, es-tabla]{babel}
\usepackage{lmodern}
\usepackage{microtype}

% ── Geometría de página ─────────────────────────────────────
\usepackage[
  top=2.8cm, bottom=3cm,
  inner=3.2cm, outer=2.5cm,
  headheight=28pt
]{geometry}

% ── Matemáticas ─────────────────────────────────────────────
\usepackage{amsmath, amssymb, amsthm}
\usepackage{mathtools}

% ── Gráficos y colores ──────────────────────────────────────
\usepackage{graphicx}
\usepackage[dvipsnames, table]{xcolor}
\usepackage{tikz}
\usetikzlibrary{arrows.meta, shapes.geometric, positioning, fit,
  backgrounds, calc, decorations.pathreplacing, shadows, matrix}

% ── Tablas ──────────────────────────────────────────────────
\usepackage{booktabs}
\usepackage{array}
\usepackage{longtable}
\usepackage{multirow}
\usepackage{tabularx}

% ── Cajas y entornos ────────────────────────────────────────
\usepackage[most]{tcolorbox}
\usepackage{mdframed}

% ── Código fuente ───────────────────────────────────────────
\usepackage{listings}

% ── Encabezados y pies ──────────────────────────────────────
\usepackage{fancyhdr}
\usepackage{emptypage}

% ── Referencias e índice ────────────────────────────────────
\usepackage{hyperref}
\usepackage{enumitem}
\usepackage{makeidx}
\makeindex

% ── Flotantes ───────────────────────────────────────────────
\usepackage{float}
\usepackage{caption}
\usepackage{subcaption}

% ── Epígrafes ───────────────────────────────────────────────
\usepackage{epigraph}
\setlength{\epigraphrule}{0pt}
\setlength{\epigraphwidth}{0.72\textwidth}

% ── Formato de títulos ──────────────────────────────────────
\usepackage{titlesec}
\usepackage{titletoc}

% ============================================================
%  PALETA DE COLORES
% ============================================================
\definecolor{TamarinBlue}  {RGB}{25,  85, 160}
\definecolor{TamarinDark}  {RGB}{10,  40,  80}
\definecolor{TamarinLight} {RGB}{210, 228, 255}
\definecolor{TamarinCyan}  {RGB}{  0, 180, 210}
\definecolor{AccentGreen}  {RGB}{ 34, 139,  34}
\definecolor{AccentOrange} {RGB}{200,  90,   0}
\definecolor{AccentRed}    {RGB}{170,  30,  20}
\definecolor{CodeBg}       {RGB}{ 22,  27,  45}
\definecolor{CodeFg}       {RGB}{215, 215, 195}
\definecolor{CodeKeyword}  {RGB}{  0, 200, 230}
\definecolor{CodeComment}  {RGB}{100, 200, 100}
\definecolor{CodeString}   {RGB}{255, 210,  80}
\definecolor{CodeNumber}   {RGB}{200, 140, 255}
\definecolor{GrisClaro}    {RGB}{245, 247, 252}
\definecolor{AmarilloClaro}{RGB}{255, 253, 220}
\definecolor{VerdeSuave}   {RGB}{230, 248, 232}
\definecolor{RojoSuave}    {RGB}{255, 235, 235}

% ============================================================
%  ESTILO DE PÁGINA
% ============================================================
\pagestyle{fancy}
\fancyhf{}
\fancyhead[LE]{\small\textcolor{TamarinBlue}{\textbf{\leftmark}}}
\fancyhead[RO]{\small\textcolor{TamarinBlue}{\textbf{\rightmark}}}
\fancyfoot[LE,RO]{\small\thepage}
\fancyfoot[C]{\small\textcolor{gray}{\textit{Guía Tamarin Prover --- UMSA 2026}}}
\renewcommand{\headrulewidth}{0.5pt}
\renewcommand{\footrulewidth}{0.3pt}

% ============================================================
%  FORMATO DE CAPÍTULOS
% ============================================================
\titleformat{\chapter}[display]
  {\normalfont\huge\bfseries\color{TamarinDark}}
  {\filleft\Large\color{TamarinBlue}
   \MakeUppercase{\chaptertitlename}\ \Huge\thechapter}
  {1ex}
  {\titlerule[1.5pt]\vspace{1ex}\filright}
  [\vspace{1ex}\titlerule]

\titleformat{\section}
  {\normalfont\Large\bfseries\color{TamarinBlue}}
  {\thesection}{1em}{}

\titleformat{\subsection}
  {\normalfont\large\bfseries\color{TamarinDark}}
  {\thesubsection}{1em}{}

\titleformat{\subsubsection}
  {\normalfont\normalsize\bfseries\color{TamarinBlue!80!black}}
  {\thesubsubsection}{1em}{}

% ============================================================
%  LENGUAJE TAMARIN
% ============================================================
\lstdefinelanguage{Tamarin}{
  morekeywords=[1]{theory, begin, end, builtins, functions, equations,
    restriction, rule, lemma, axiom, let, in, Fr, In, Out,
    All, Ex, not, and, or, exists-trace, all-traces},
  morekeywords=[2]{K, KU, KD, true, false},
  morekeywords=[3]{sign, verify, pk, senc, sdec, encaps, decaps,
    ek_from_dk, h, aenc, adec},
  morestring=[b]{"},
  morestring=[b]{'},
  morecomment=[l]{//},
  morecomment=[s]{/*}{*/},
  sensitive=true,
}

\lstdefinestyle{TamarinStyle}{
  language=Tamarin,
  backgroundcolor=\color{CodeBg},
  basicstyle=\ttfamily\small\color{CodeFg},
  keywordstyle=[1]{\color{CodeKeyword}\bfseries},
  keywordstyle=[2]{\color{CodeNumber}\bfseries},
  keywordstyle=[3]{\color{AccentOrange!90!white}},
  stringstyle=\color{CodeString},
  commentstyle=\color{CodeComment}\itshape,
  numbers=left,
  numberstyle=\tiny\color{gray!60},
  numbersep=10pt,
  breaklines=true,
  frame=none,
  tabsize=4,
  keepspaces=true,
  xleftmargin=16pt,
  xrightmargin=4pt,
  showstringspaces=false,
  captionpos=b,
}

\lstdefinestyle{BashStyle}{
  language=bash,
  backgroundcolor=\color{TamarinDark},
  basicstyle=\ttfamily\small\color{white},
  commentstyle=\color{CodeComment}\itshape,
  stringstyle=\color{CodeString},
  numbers=none,
  frame=none,
  breaklines=true,
  tabsize=2,
  xleftmargin=10pt,
  showstringspaces=false,
}

\lstset{style=TamarinStyle}

% ============================================================
%  CAJAS TEMÁTICAS
% ============================================================
\tcbuselibrary{most}

\newtcolorbox{cajaDef}[2][]{
  colback=TamarinLight!50, colframe=TamarinBlue, coltitle=white,
  fonttitle=\bfseries\small, title={#2}, #1,
  boxrule=1pt, arc=6pt, left=10pt, right=10pt, top=6pt, bottom=6pt}

\newtcolorbox{cajaEjemplo}[2][]{
  colback=VerdeSuave, colframe=AccentGreen, coltitle=white,
  fonttitle=\bfseries\small, title={Ejemplo: #2}, #1,
  boxrule=1pt, arc=6pt, left=10pt, right=10pt}

\newtcolorbox{cajaAnalogia}[2][]{
  colback=AmarilloClaro, colframe=AccentOrange, coltitle=white,
  fonttitle=\bfseries\small, title={Analog\'ia: #2}, #1,
  boxrule=1pt, arc=6pt, left=10pt, right=10pt}

\newtcolorbox{cajaAviso}[2][]{
  colback=RojoSuave, colframe=AccentRed, coltitle=white,
  fonttitle=\bfseries\small, title={Importante: #2}, #1,
  boxrule=1pt, arc=6pt, left=10pt, right=10pt}

\newtcolorbox{cajaResultado}[2][]{
  colback=VerdeSuave, colframe=AccentGreen!80!black, coltitle=white,
  fonttitle=\bfseries\normalsize, title={\textcolor{white}{#2}}, #1,
  boxrule=1.5pt, arc=8pt, left=10pt, right=10pt}

\newtcolorbox{cajaSintaxis}[2][]{
  colback=CodeBg!95, colframe=TamarinCyan, coltitle=TamarinCyan,
  fonttitle=\bfseries\small\ttfamily, title={Sintaxis: #2}, #1,
  boxrule=1pt, arc=6pt, left=10pt, right=10pt,
  fontupper=\ttfamily\small\color{CodeFg}}

\definecolor{LogicBlue}    {RGB}{ 20,  45, 110}
\definecolor{LogicBg}      {RGB}{243, 246, 255}

\newtcolorbox{cajaResumen}[1][]{
  colback=TamarinLight, colframe=TamarinBlue,
  title={\bfseries\color{TamarinDark} Resumen del Cap\'itulo}, #1,
  boxrule=1.2pt, arc=8pt, left=10pt, right=10pt}

\newtcolorbox{cajaLogica}[2][]{
  colback=LogicBg, colframe=LogicBlue, coltitle=white,
  fonttitle=\bfseries\small, title={Formulaci\'on en L\'ogica Formal: #2}, #1,
  boxrule=1.2pt, arc=6pt, left=10pt, right=10pt, top=7pt, bottom=7pt}

\newtheorem{teorema}{Teorema}[chapter]
\newtheorem{definicion}{Definici\'on}[chapter]

% ============================================================
%  HYPERREF
% ============================================================
\hypersetup{
  colorlinks=true, linkcolor=TamarinBlue, urlcolor=TamarinCyan,
  citecolor=AccentGreen,
  pdftitle={Gu\'ia Tamarin Prover --- Protocolo PQC UMSA 2026},
  pdfauthor={UMSA --- Carrera de Inform\'atica --- 2026},
  pdfsubject={Verificaci\'on formal, Tamarin, criptograf\'ia post-cu\'antica},
  bookmarksopen=true, bookmarksnumbered=true,
}

\setcounter{secnumdepth}{3}
\setcounter{tocdepth}{2}

% ============================================================
\begin{document}
% ============================================================

% ── PORTADA ─────────────────────────────────────────────────
\begin{titlepage}
  \pagecolor{TamarinDark}
  \color{white}
  \centering
  \vspace*{1.5cm}

  \begin{tikzpicture}
    \node[draw=TamarinCyan, line width=2pt, rounded corners=10pt,
          inner sep=16pt, fill=TamarinBlue!40] {
      \Large\bfseries\color{white}
      \begin{tabular}{c}
        UMSA\\
        Universidad Mayor de San Andr\'es\\
        {\normalsize Carrera de Inform\'atica --- FCPN}
      \end{tabular}
    };
  \end{tikzpicture}

  \vspace{2cm}

  {\Huge\bfseries\color{TamarinCyan}
    Gu\'ia de \\ Tamarin Prover}\\[0.8em]
  {\Large\color{white!85}
    Verificaci\'on Formal de Protocolos Criptogr\'aficos}\\[0.4em]
  {\large\color{white!70}
    Fundamentos, Sintaxis y Aplicaci\'on al\\
    Protocolo de Votaci\'on Electr\'onica Post-Cu\'antica}

  \vspace{1.5cm}

  \begin{tcolorbox}[
    colback=TamarinBlue!60, colframe=TamarinCyan,
    width=0.85\textwidth, arc=8pt, boxrule=1.5pt
  ]
    \centering\itshape\large\color{white}
    ``La verificaci\'on formal no prueba que un protocolo\\
    funciona bien en algunos casos ---\\
    prueba que \textbf{no puede fallar} en ning\'un caso posible.''
  \end{tcolorbox}

  \vfill

  \begin{tikzpicture}
    \node[fill=TamarinBlue!30, rounded corners=6pt,
          inner sep=12pt, text width=0.7\textwidth, align=center] {
      {\bfseries\color{TamarinCyan} Proyecto de Grado 2026}\\[4pt]
      {\small\color{white!80}
        ML-KEM-768 $\cdot$ ML-DSA-65 $\cdot$ FIPS 203/204\\
        Protocolo: \texttt{tamarin/protocolo.spthy}\\
        Motor: Maude 3.5.1 $+$ Tamarin Prover}
    };
  \end{tikzpicture}

  \vspace{1cm}
  \nopagecolor
\end{titlepage}

\newpage\thispagestyle{empty}\null\newpage

% ── DEDICATORIA ─────────────────────────────────────────────
\thispagestyle{empty}
\vspace*{5cm}
\begin{flushright}
  \itshape\large
  A todos los que creen que\\
  la matem\'atica puede construir\\
  sistemas m\'as justos y confiables.
\end{flushright}
\newpage

% ── PRÓLOGO ─────────────────────────────────────────────────
\chapter*{Pr\'ologo}
\addcontentsline{toc}{chapter}{Pr\'ologo}

Esta gu\'ia naci\'o de la necesidad de explicar, con rigor pero
tambi\'en con claridad, c\'omo se usa \textbf{Tamarin Prover} para
demostrar formalmente que un protocolo criptogr\'afico es seguro.
Fue escrita como parte del desarrollo de la tesis de grado sobre un
\textit{protocolo de votaci\'on electr\'onica post-cu\'antica},
dise\~nado para resistir ataques de computadoras cu\'anticas mediante
los est\'andares FIPS~203 (ML-KEM-768) y FIPS~204 (ML-DSA-65) del NIST.

\medskip

El lector encontrar\'a aqu\'i tanto los fundamentos te\'oricos de la
verificaci\'on formal como una gu\'ia pr\'actica, paso a paso, del
lenguaje \texttt{.spthy} de Tamarin. Los ejemplos de c\'odigo est\'an
tomados directamente del archivo \texttt{protocolo.spthy} que acompa\~na
la tesis, de modo que cada concepto abstracto tiene un correlato
concreto y ejecutable.

\medskip

\noindent La gu\'ia est\'a organizada en \textbf{ocho cap\'itulos}:
los primeros tres cubren los fundamentos conceptuales; los cap\'itulos
cuatro a seis exploran la sintaxis y sem\'antica de Tamarin; el
cap\'itulo siete presenta el an\'alisis completo del protocolo PQC;
y el cap\'itulo ocho cierra con las instrucciones de ejecuci\'on y
los resultados obtenidos.

\medskip\noindent
\textit{La Paz, Bolivia, 2026}

% ── TABLA DE CONTENIDOS ─────────────────────────────────────
\tableofcontents
\listoffigures
\listoftables

% ────────────────────────────────────────────────────────────
%   PARTE I — FUNDAMENTOS
% ────────────────────────────────────────────────────────────
\part{Fundamentos de la Verificaci\'on Formal}

% ============================================================
\chapter{Por qu\'e Verificar Formalmente un Protocolo}
% ============================================================

\epigraph{%
  \textit{``Anyone can invent a cryptographic system that they
  themselves cannot break. What's hard is creating one that
  others cannot break.''}
}{--- Bruce Schneier}

\section{El problema de los protocolos ``casi seguros''}

La historia de la criptograf\'ia moderna es, en parte, un cat\'alogo
de protocolos que \textit{parec\'ian} seguros pero no lo eran.
El protocolo WEP, ampliamente utilizado en la d\'ecada de 2000,
fue dise\~nado por ingenieros expertos y tard\'o a\~nos en ser
completamente roto. Lo mismo ocurri\'o con SSL~2.0, con MS-CHAP,
y con docenas de esquemas que resistieron pruebas emp\'iricas durante
a\~nos antes de ser vulnerados.

\medskip

El problema fundamental es que probar un sistema con casos de prueba,
por exhaustivo que sea el proceso, \textbf{nunca puede garantizar}
la ausencia de vulnerabilidades: solo prueba que el sistema funciona
en los casos probados. Un atacante inteligente no ataca donde usted
prob\'o; ataca donde usted no prob\'o.

\begin{cajaAnalogia}{El duelo de espadas}
  Probar un protocolo con pruebas unitarias es como entrenar a un
  espadach\'in contra rivales conocidos. Por m\'as que practique
  con cien contrincantes distintos, siempre existe un movimiento
  que ninguno de ellos us\'o. La verificaci\'on formal es distinta:
  en lugar de entrenar contra rivales concretos, se demuestra
  matem\'aticamente que \textit{no existe ning\'un movimiento posible}
  que pueda derrotar al espadach\'in, sin importar qui\'en lo enfrente.
\end{cajaAnalogia}

\section{Qu\'e garantiza la verificaci\'on formal}

\begin{cajaDef}{Verificaci\'on Formal de Protocolos}
  La \textbf{verificaci\'on formal} de un protocolo criptogr\'afico
  es un proceso matem\'atico que demuestra, con certeza absoluta dentro
  de un modelo declarado, que el protocolo cumple ciertas propiedades
  de seguridad bajo \emph{todos} los posibles comportamientos del
  ambiente y del atacante, sin excepci\'on.
\end{cajaDef}

\medskip

Las propiedades que t\'ipicamente se verifican son:
\begin{itemize}[leftmargin=*, itemsep=4pt]
  \item \textbf{Confidencialidad}: un dato secreto nunca es revelado
    al atacante, sin importar qu\'e mensajes intercepte o inyecte.
  \item \textbf{Autenticidad}: el receptor de un mensaje puede estar
    seguro de que fue enviado por el emisor leg\'itimo.
  \item \textbf{Integridad}: el contenido de un mensaje no puede ser
    alterado en tr\'ansito sin ser detectado.
  \item \textbf{No repudio}: el emisor no puede negar haber enviado
    un mensaje que firm\'o.
  \item \textbf{Unicidad / anti-replay}: un mensaje v\'alido no puede
    ser reutilizado por un atacante para obtener efectos adicionales.
\end{itemize}

\section{Herramientas de verificaci\'on formal}

\begin{table}[H]
\centering
\caption{Comparaci\'on de herramientas de verificaci\'on formal.}
\label{tab:herramientas}
\begin{tabularx}{\textwidth}{>{\bfseries}p{3.2cm} X X}
\toprule
Herramienta & Enfoque & Caracter\'isticas principales \\
\midrule
\textbf{Tamarin} & Reescritura de multiconjuntos &
  Motor Maude; estado mutable, XOR, DH; altamente expresivo.\\[4pt]
ProVerif & \'Algebra de procesos &
  M\'as autom\'atico; menos expresivo para equivalencias.\\[4pt]
Scyther & Grafos de roles &
  M\'as sencillo; menos potente para protocolos con estado.\\[4pt]
CPSA & \'Algebra de hebras &
  Especializado en intercambio de claves.\\
\bottomrule
\end{tabularx}
\end{table}

\medskip

En este trabajo usamos \textbf{Tamarin Prover} porque:
\begin{enumerate}[itemsep=2pt]
  \item Soporta el modelo de estado global mutable, necesario para
    modelar el \'anfora digital.
  \item Permite definir restricciones arbitrarias como axiomas del
    sistema (por ejemplo, ``cada votante solo puede votar una vez'').
  \item Maneja la ecuaci\'on de correcci\'on del KEM
    \texttt{decaps(dk, encaps(ek, r)) = r} sin necesidad de plugins.
\end{enumerate}

\begin{cajaResumen}
  La verificaci\'on formal demuestra la seguridad de un protocolo
  bajo \emph{cualquier} comportamiento del atacante, no solo bajo
  los ataques conocidos. Tamarin Prover es la herramienta elegida
  por su capacidad para manejar estado mutable y ecuaciones
  criptogr\'aficas arbitrarias.
\end{cajaResumen}

% ============================================================
\chapter{El Modelo de Atacante Dolev-Yao}
% ============================================================

\epigraph{%
  \textit{``The attacker knows everything about the protocol, controls
  the network completely, and is bounded only by the laws of
  mathematics.''}
}{--- Adaptado de Dolev \& Yao, 1983}

\section{Origen y motivaci\'on}

En 1983, Danny Dolev y Andrew Yao formalizaron por primera vez un
modelo matem\'atico del adversario en protocolos criptogr\'aficos:

\begin{cajaDef}{Modelo de Atacante Dolev-Yao}
  En el modelo \textbf{Dolev-Yao}, el atacante $\mathcal{A}$ controla
  completamente el canal de red. Puede:
  \begin{itemize}[itemsep=2pt]
    \item \textbf{Interceptar} cualquier mensaje enviado por cualquier
      participante.
    \item \textbf{Modificar} mensajes en tr\'ansito.
    \item \textbf{Eliminar} mensajes (denegaci\'on de servicio).
    \item \textbf{Reenviar} mensajes previos (ataques de repetici\'on).
    \item \textbf{Construir} nuevos mensajes a partir de fragmentos
      que ha observado.
    \item \textbf{Hacerse pasar} por cualquier participante leg\'itimo.
  \end{itemize}
  Sin embargo, $\mathcal{A}$ \emph{no puede} romper las primitivas
  criptogr\'aficas: no puede descifrar sin la clave, no puede falsificar
  firmas, y no puede invertir funciones hash.
\end{cajaDef}

\section{El atacante en Tamarin: el hecho \texttt{K}}

En Tamarin, el conocimiento del atacante se representa con el hecho
especial \texttt{K(t)}, que significa ``el atacante conoce el
t\'ermino~$t$''.

\begin{cajaSintaxis}{Interacci\'on con la red}
  \texttt{Out(t)} $\rightarrow$ el protocolo env\'ia \textit{t} a la red\\
  \texttt{In(t)} $\rightarrow$ el protocolo recibe \textit{t} de la red\\
  \texttt{K(t)} $\rightarrow$ el atacante conoce el t\'ermino \textit{t}
\end{cajaSintaxis}

\begin{cajaAviso}{El atacante controla \texttt{In}}
  Cuando una regla tiene \texttt{In(t)} en sus premisas, \textbf{no
  significa que el mensaje lleg\'o del participante leg\'itimo}.
  Significa que lleg\'o de la red, bajo control del atacante.
  El atacante puede haber modificado, sustituido o fabricado ese mensaje.
\end{cajaAviso}

\section{Por qu\'e Dolev-Yao en un sistema electoral}

En el contexto de una votaci\'on electr\'onica, el modelo Dolev-Yao
captura amenazas muy reales:
\begin{itemize}[leftmargin=*, itemsep=4pt]
  \item Un atacante que compromete un enrutador y modifica el tr\'afico.
  \item Un administrador de red que intenta inyectar boletas falsas.
  \item Un atacante que graba tr\'afico y lo reenv\'ia (replay).
  \item Un gobierno extranjero con capacidad de interceptar
    comunicaciones.
\end{itemize}

\begin{cajaResumen}
  El modelo Dolev-Yao pone al atacante m\'as poderoso posible en el
  control de la red. Si el protocolo es seguro en este modelo, es
  seguro contra cualquier ataque de red concebible. El hecho
  \texttt{K(t)} en Tamarin representa el conocimiento acumulado
  del atacante a lo largo de la ejecuci\'on del protocolo.
\end{cajaResumen}

% ============================================================
\chapter{Criptograf\'ia Simb\'olica vs.\ Computacional}
% ============================================================

\epigraph{%
  \textit{``Symbolic models trade precision for automation;
  computational models trade automation for precision.''}
}{--- Cortier, Kremer \& Warinschi, 2011}

\section{Dos enfoques para modelar la criptograf\'ia}

\begin{table}[H]
\centering
\caption{Comparaci\'on entre modelos simb\'olico y computacional.}
\label{tab:simbol-comput}
\begin{tabularx}{\textwidth}{>{\bfseries}p{3.5cm} X X}
\toprule
Aspecto & Modelo Simb\'olico & Modelo Computacional \\
\midrule
Representaci\'on & T\'erminos algebraicos & Cadenas de bits \\[3pt]
Operaciones & Funciones opacas & Algoritmos reales con costos\\[3pt]
Atacante & Dolev-Yao & Adversario probabil\'istico\\[3pt]
Garant\'ia & Seguridad en el modelo formal & Seguridad computacional\\[3pt]
Automatizaci\'on & Alta (Tamarin, ProVerif) & Baja (pruebas manuales)\\[3pt]
Aplicaci\'on & Dise\~no de protocolos & Criptograf\'ia de bajo nivel\\
\bottomrule
\end{tabularx}
\end{table}

\section{El modelo simb\'olico: cajas negras perfectas}

\begin{cajaDef}{Caja Negra Perfecta}
  Una operaci\'on criptogr\'afica es una \textbf{caja negra perfecta}
  en el modelo simb\'olico si:
  \begin{enumerate}[itemsep=2pt]
    \item Solo se puede invertir con el conocimiento correcto (la clave).
    \item No existe ning\'un atajo computacional.
    \item El atacante no puede aprender nada sobre las entradas a
      partir de la salida, excepto lo que las ecuaciones del modelo
      permiten.
  \end{enumerate}
\end{cajaDef}

\begin{lstlisting}[caption={Declaraci\'on de primitivas como cajas negras en Tamarin.}]
builtins: dest-signing, symmetric-encryption

functions:
    encaps/2,        // ML-KEM-768: encapsula un secreto k con pk
    decaps/2,        // ML-KEM-768: desencapsula con dk, recupera k
    ek_from_dk/1,    // Deriva la clave publica de encapsulamiento
    h/1              // SHA3-256: hash para recibos de voto

equations:
    // Propiedad de correccion del KEM (FIPS 203)
    decaps(dk, encaps(ek_from_dk(dk), r)) = r
\end{lstlisting}

\begin{table}[H]
\centering
\caption{Builtins de Tamarin Prover y sus funciones.}
\label{tab:builtins}
\begin{tabular}{>{\ttfamily}p{5cm} p{8cm}}
\toprule
\normalfont Builtin & Funciones que incluye \\
\midrule
dest-signing & \texttt{sign/2}, \texttt{verify/3}, \texttt{pk/1} \\[3pt]
symmetric-encryption & \texttt{senc/2}, \texttt{sdec/2} \\[3pt]
asymmetric-encryption & \texttt{aenc/2}, \texttt{adec/2}, \texttt{pk/1} \\[3pt]
hashing & \texttt{h/1} \\[3pt]
xor & Operaci\'on XOR ($\oplus$) \\[3pt]
diffie-hellman & Aritm\'etica modular para DH \\
\bottomrule
\end{tabular}
\end{table}

\section{Ecuaciones: las propiedades de las funciones}

\begin{cajaEjemplo}{Ecuaciones est\'andar de Tamarin}
  Las siguientes ecuaciones est\'an impl\'icitas en los builtins:
  \begin{itemize}[itemsep=3pt]
    \item Firma digital:
      $\texttt{verify(sign(m, sk), m, pk(sk)) = true}$
    \item Cifrado sim\'etrico:
      $\texttt{sdec(senc(m, k), k) = m}$
    \item KEM (declarada expl\'icitamente):
      $\texttt{decaps(dk, encaps(ek\_from\_dk(dk), r)) = r}$
  \end{itemize}
\end{cajaEjemplo}

\begin{cajaLogica}{Teor\'ia Ecuacional y Congruencia Simb\'olica}
  En el modelo simb\'olico de Tamarin, los mensajes son elementos del \'algebra libre de t\'erminos $\mathcal{T}_\Sigma(\mathcal{V})$ sobre una signatura funcional $\Sigma$ y variables $\mathcal{V}$, cocientada por una teor\'ia ecuacional $E$ ($\mathcal{T}_\Sigma(\mathcal{V}) / =_E$).
  
  \textbf{Proposiciones Ecuacionales del Protocolo PQC:}
  \begin{align*}
    \forall m \in \mathcal{T},\; \forall sk \in \mathcal{V}: & \quad
      \operatorname{verify}(\operatorname{sign}(m, sk),\, m,\, \operatorname{pk}(sk)) =_E \mathrm{true} \\
    \forall m, k \in \mathcal{T}: & \quad
      \operatorname{sdec}(\operatorname{senc}(m, k),\, k) =_E m \\
    \forall dk, r \in \mathcal{V}: & \quad
      \operatorname{decaps}(dk,\, \operatorname{encaps}(\operatorname{ek\_from\_dk}(dk),\, r)) =_E r
  \end{align*}
  \textbf{Axioma de Caja Negra (No-deribabilidad):} Para cualquier t\'ermino cifrado $\operatorname{senc}(m, k)$, se cumple que el conocimiento del atacante est\'a acotado por:
  \[
    \mathrm{K}(\operatorname{senc}(m, k)) \;\land\; \neg\,\mathrm{K}(k) \implies \neg\,\mathrm{K}(m)
  \]
\end{cajaLogica}

\begin{cajaAviso}{Qu\'e no pueden hacer las ecuaciones}
  Las ecuaciones solo definen las formas \emph{leg\'itimas} de usar
  las funciones. No permiten al atacante obtener informaci\'on
  adicional: a partir de \texttt{senc(m, k)}, sin conocer \texttt{k},
  el atacante no puede aprender nada sobre \texttt{m}.
\end{cajaAviso}

\begin{cajaResumen}
  El modelo simb\'olico trata las operaciones criptogr\'aficas como
  cajas negras perfectas. Las ecuaciones capturan \'unicamente las
  propiedades de correcci\'on. Este enfoque permite la automatizaci\'on
  completa de la verificaci\'on.
\end{cajaResumen}

% ────────────────────────────────────────────────────────────
%   PARTE II — TAMARIN: LENGUAJE Y SEMÁNTICA
% ────────────────────────────────────────────────────────────
\part{Tamarin Prover: Lenguaje y Sem\'antica}

% ============================================================
\chapter{Arquitectura y Estructura de un Archivo \texttt{.spthy}}
% ============================================================

\epigraph{%
  \textit{``A formal model is a contract: it defines exactly what
  we are proving and exactly under what assumptions.''}
}{--- Cas Cremers, ETH Z\"urich}

\section{Estructura general de una teor\'ia Tamarin}

Un archivo \texttt{.spthy} (\textit{security protocol theory}) es
la unidad b\'asica de especificaci\'on en Tamarin:

\begin{lstlisting}[caption={Esqueleto de una teor\'ia Tamarin.}]
theory NombreDeTeoria
begin
  // 1. Primitivas (builtins y funciones)
  builtins: ...
  functions: ...
  equations: ...

  // 2. Restricciones (axiomas del sistema)
  restriction NombreRestriccion: "formula"

  // 3. Reglas del protocolo
  rule NombreRegla:
      [ Premisas ]
    --[ Acciones ]->
      [ Conclusiones ]

  // 4. Lemas de seguridad
  lemma NombreLema: tipo "formula"
end
\end{lstlisting}

\begin{cajaDef}{Teor\'ia Tamarin}
  Una \textbf{teor\'ia} en Tamarin es un m\'odulo autocontenido que
  especifica:
  \begin{itemize}[itemsep=2pt]
    \item Las primitivas criptogr\'aficas disponibles.
    \item Las reglas de reescritura que definen el comportamiento
      del protocolo y del atacante.
    \item Las restricciones que acotan el espacio de trazas v\'alidas.
    \item Los lemas que Tamarin debe verificar (probar o refutar).
  \end{itemize}
\end{cajaDef}

\section{El concepto de multiconjunto de hechos}

El estado del sistema en Tamarin se representa como un
\textbf{multiconjunto de hechos} (\textit{multiset of facts}).

\begin{cajaSintaxis}{Hechos en Tamarin}
  \texttt{NombreHecho(arg1, arg2, \ldots)}\\
  \texttt{!} prefijo $\Rightarrow$ hecho \textbf{persistente}
    (no se consume al usarse)\\
  Sin \texttt{!} $\Rightarrow$ hecho \textbf{lineal}
    (se consume una vez)
\end{cajaSintaxis}

\begin{figure}[H]
\centering
\begin{tikzpicture}[
  box/.style={draw=TamarinBlue, fill=TamarinLight!70, rounded corners=5pt,
              inner sep=8pt, text width=3.8cm, align=center, font=\small},
  arrow/.style={-Stealth, thick, TamarinBlue},
  lbl/.style={font=\small\itshape, color=gray, align=center}
]
  \node[box] (e1) {\textbf{Estado inicial}\\[4pt]
     \texttt{Fr(\textasciitilde sk\_AR)}\\
     \texttt{Fr(\textasciitilde sk\_SV)}};
  \node[box, right=2.8cm of e1] (e2) {\textbf{Tras Setup\_AR}\\[4pt]
     \texttt{!ClavePrivadaAR}\\
     \texttt{!ClavePublicaAR}};
  \node[box, right=2.8cm of e2] (e3) {\textbf{Tras Setup\_SV}\\[4pt]
     \texttt{!ClaveDesenc..}\\
     \texttt{!ClavePublicaKEM..}};
  \draw[arrow] (e1) -- node[above, lbl]{regla\\Setup\_AR} (e2);
  \draw[arrow] (e2) -- node[above, lbl]{regla\\Setup\_SV} (e3);
\end{tikzpicture}
\caption{Evoluci\'on del multiconjunto de hechos durante el protocolo.}
\label{fig:multiconjunto}
\end{figure}

\section{T\'erminos y variables}

En Tamarin, los \textbf{t\'erminos} son las ``palabras'' del lenguaje:
\begin{itemize}[leftmargin=*, itemsep=4pt]
  \item \textbf{Variables frescas} (\texttt{\textasciitilde x}):
    valores aleatorios \'unicos generados en tiempo de ejecuci\'on.
  \item \textbf{Variables de patr\'on} (\texttt{x}):
    variables que se instancian al hacer coincidir patrones.
  \item \textbf{Constantes p\'ublicas} (\texttt{'nombre'}):
    cadenas literales conocidas por todos, incluyendo el atacante.
  \item \textbf{Tuplas} (\texttt{<a, b, c>}): agrupaciones de t\'erminos.
  \item \textbf{Aplicaciones de funci\'on} (\texttt{f(a, b)}).
\end{itemize}

\begin{cajaEjemplo}{Variables en \texttt{protocolo.spthy}}
  \begin{itemize}[itemsep=3pt]
    \item \texttt{\textasciitilde sk\_AR}: clave privada fresca de la AR.
    \item \texttt{\textasciitilde k\_efimera}: secreto KEM ef\'imero.
    \item \texttt{'AR'}: constante p\'ublica que identifica el rol AR.
    \item \texttt{<idV, pk\_V, ts\_reg>}: tupla de la credencial.
  \end{itemize}
\end{cajaEjemplo}

\section{Trazas de ejecuci\'on}

\begin{cajaDef}{Traza de Ejecuci\'on}
  Una \textbf{traza} $\tau = \langle (S_0, A_0), (S_1, A_1), \ldots,
  (S_n, A_n) \rangle$ es una secuencia donde:
  \begin{itemize}[itemsep=2pt]
    \item $S_i$ es el multiconjunto de hechos en el paso $i$.
    \item $A_i$ es el multiconjunto de acciones realizadas.
    \item Cada transici\'on $S_i \to S_{i+1}$ est\'a justificada por
      la aplicaci\'on de alguna regla del protocolo.
  \end{itemize}
\end{cajaDef}

\begin{cajaResumen}
  Un archivo \texttt{.spthy} define una \emph{teor\'ia}: las
  primitivas, las reglas de reescritura, las restricciones y los
  lemas. El estado del protocolo es un multiconjunto de hechos que
  evoluciona por aplicaci\'on de reglas. Las trazas son las historias
  de ejecuci\'on posibles del sistema.
\end{cajaResumen}

% ============================================================
\chapter{Reglas: El Coraz\'on del Modelo}
% ============================================================

\epigraph{%
  \textit{``Rules are the atomic steps of the protocol, each one
  describing a transition from one global state to another.''}
}{--- Manual de Tamarin Prover, 2023}

\section{Sintaxis de una regla}

\begin{cajaSintaxis}{Regla de reescritura}
  \texttt{rule NombreRegla:}\\
  \texttt{\ \ \ \ let var = expresion in}\\
  \texttt{\ \ \ \ [ Premisas ] --[ Acciones ]-> [ Conclusiones ]}
\end{cajaSintaxis}

\begin{itemize}[leftmargin=*, itemsep=6pt]
  \item \textbf{Premisas}: hechos presentes en el multiconjunto de estado.
    Los lineales (sin \texttt{!}) se consumen; los persistentes no.
  \item \textbf{Acciones}: etiquetas registradas en la traza de ejecuci\'on.
    Son los puntos de observaci\'on que usan los lemas.
  \item \textbf{Conclusiones}: hechos nuevos que se agregan al estado.
\end{itemize}

\section{Hechos especiales: Fr, In, Out}

\begin{table}[H]
\centering
\caption{Hechos especiales de Tamarin y su significado.}
\label{tab:hechos-esp}
\begin{tabular}{>{\ttfamily}p{3.2cm} p{9.5cm}}
\toprule
\normalfont Hecho & Significado \\
\midrule
Fr(\textasciitilde x) &
  Genera un valor \textbf{fresco} (aleatorio, nunca antes visto).
  Solo puede aparecer en premisas. \\[4pt]
In(t) &
  La regla \textbf{recibe} el t\'ermino \texttt{t} de la red.
  El atacante controla qu\'e valor se instancia aqu\'i. \\[4pt]
Out(t) &
  La regla \textbf{env\'ia} el t\'ermino \texttt{t} a la red.
  El atacante puede leer y registrar este valor. \\
\bottomrule
\end{tabular}
\end{table}

\section{Hechos lineales vs.\ persistentes}

\begin{figure}[H]
\centering
\begin{tikzpicture}[
  box/.style={rounded corners=5pt, inner sep=10pt,
              text width=5cm, align=center, font=\small},
  arrow/.style={-Stealth, thick}
]
  \node[box, fill=AccentOrange!20, draw=AccentOrange] (lineal) at (0,0) {
    \textbf{Hecho Lineal}\\[4pt]
    \texttt{AnforaDigital(...)}\\[4pt]
    Se \emph{consume} cuando una regla lo usa como premisa.
    Puede usarse \textbf{una sola vez}.
  };
  \node[box, fill=AccentGreen!15, draw=AccentGreen] (persist) at (8,0) {
    \textbf{Hecho Persistente}\\[4pt]
    \texttt{!ClavePrivadaAR(...)}\\[4pt]
    \emph{No se consume} cuando una regla lo usa.
    Puede usarse \textbf{m\'ultiples veces}.
  };
  \node[below=0.6cm of lineal, font=\small\itshape, text=AccentOrange]
    {El \'anfora se ``abre'' exactamente una vez};
  \node[below=0.6cm of persist, font=\small\itshape, text=AccentGreen]
    {La AR firma m\'ultiples credenciales};
\end{tikzpicture}
\caption{Diferencia entre hechos lineales y persistentes.}
\label{fig:lineal-persist}
\end{figure}

\section{An\'alisis de las seis reglas del protocolo PQC}

\subsection{Regla 1: Setup\_Autoridad\_Registro}

\begin{lstlisting}[caption={Configuraci\'on de la Autoridad de Registro con ML-DSA-65.}]
rule Setup_Autoridad_Registro:
    let  pk_AR = pk(~sk_AR)  in
    [ Fr(~sk_AR) ]
  --[ SetupAR(pk_AR) ]->
    [ !ClavePrivadaAR('AR', ~sk_AR)   // Confidencial
    , !ClavePublicaAR('AR', pk_AR)
    , Out(pk_AR)                       // Publicada en la red
    ]
\end{lstlisting}

\begin{cajaDef}{Lectura en lenguaje llano}
  Si est\'a disponible un valor aleatorio fresco \texttt{\textasciitilde
  sk\_AR}, entonces se establece como clave privada de la AR
  (persistente, confidencial), se registra la clave p\'ublica
  correspondiente (persistente), y se publica \texttt{pk\_AR} en la
  red (cualquiera puede verla).
\end{cajaDef}

\begin{cajaLogica}{Sem\'antica MSR y L\'ogica Lineal de Setup\_Autoridad\_Registro}
  \textbf{Transici\'on de Estado en Reescritura de Multiconjuntos (MSR):}
  \[
    \begin{aligned}
      \mathcal{S} \cup \{ \operatorname{Fr}(\mathit{sk}_{\mathrm{AR}}) \}
      &\xrightarrow{\;\operatorname{SetupAR}(\operatorname{pk}(\mathit{sk}_{\mathrm{AR}}))\;} \\
      \mathcal{S} \cup \big\{
        &!\operatorname{ClavePrivadaAR}(\text{'AR'}, \mathit{sk}_{\mathrm{AR}}), \\
        &!\operatorname{ClavePublicaAR}(\text{'AR'}, \operatorname{pk}(\mathit{sk}_{\mathrm{AR}})),\,
        \operatorname{Out}(\operatorname{pk}(\mathit{sk}_{\mathrm{AR}}))
      \big\}
    \end{aligned}
  \]
  \textbf{Proposici\'on en L\'ogica Lineal de Girard:}
  \[
    \begin{aligned}
      \operatorname{Fr}(\mathit{sk}_{\mathrm{AR}}) \multimap
      &\operatorname{SetupAR}(\operatorname{pk}(\mathit{sk}_{\mathrm{AR}})) \;\otimes\;
      !\operatorname{ClavePrivadaAR}(\text{'AR'}, \mathit{sk}_{\mathrm{AR}}) \\
      &\otimes\; !\operatorname{ClavePublicaAR}(\text{'AR'}, \operatorname{pk}(\mathit{sk}_{\mathrm{AR}})) \;\otimes\;
      \operatorname{Out}(\operatorname{pk}(\mathit{sk}_{\mathrm{AR}}))
    \end{aligned}
  \]
  \textbf{Condici\'on L\'ogica de Traza Temporal (FOTL):}
  \[
    \forall \mathit{pk}, t.\; \operatorname{SetupAR}(\mathit{pk})@t \implies
    \exists \mathit{sk}.\; \big(\mathit{pk} = \operatorname{pk}(\mathit{sk}) \;\land\;
    \exists t' \le t.\; \operatorname{Fr}(\mathit{sk})@t'\big)
  \]
\end{cajaLogica}

\subsection{Regla 2: Setup\_Servidor\_Votacion}

\begin{lstlisting}[caption={SV genera claves ML-KEM-768 y ML-DSA-65.}]
rule Setup_Servidor_Votacion:
    let
        ek_SV = ek_from_dk(~dk_SV)   // Clave publica KEM
        pk_SV = pk(~sk_SV)           // Clave publica firma
    in
    [ Fr(~dk_SV), Fr(~sk_SV) ]
  --[ SetupSV(ek_SV, pk_SV) ]->
    [ !ClaveDesencapsulamientoAC('AC', ~dk_SV) // Solo la AC
    , !ClavePublicaKEM_SV('SV', ek_SV)
    , !ClavePublicaFirma_SV('SV', pk_SV)
    , Out(ek_SV)
    , Out(pk_SV)
    ]
\end{lstlisting}

\textbf{Punto clave:} \texttt{\textasciitilde dk\_SV} nunca aparece en
un hecho \texttt{Out}; solo existe en el estado de la \texttt{'AC'}.
Esto modela la \textit{separaci\'on de custodia}: el SV no puede
descifrar los votos que custodia.

\begin{cajaLogica}{Sem\'antica MSR y L\'ogica Lineal de Setup\_Servidor\_Votacion}
  \textbf{Transici\'on de Estado MSR:}
  \[
    \begin{aligned}
      \mathcal{S} \cup \{ \operatorname{Fr}(\mathit{dk}_{\mathrm{SV}}), \operatorname{Fr}(\mathit{sk}_{\mathrm{SV}}) \}
      &\xrightarrow{\;\operatorname{SetupSV}(\mathit{ek}_{\mathrm{SV}}, \mathit{pk}_{\mathrm{SV}})\;} \\
      \mathcal{S} \cup \big\{
        &!\operatorname{ClaveDesencapsulamientoAC}(\text{'AC'}, \mathit{dk}_{\mathrm{SV}}),\,
        !\operatorname{ClavePublicaKEM\_SV}(\text{'SV'}, \mathit{ek}_{\mathrm{SV}}), \\
        &!\operatorname{ClavePublicaFirma\_SV}(\text{'SV'}, \mathit{pk}_{\mathrm{SV}}),\,
        \operatorname{Out}(\mathit{ek}_{\mathrm{SV}}),\,
        \operatorname{Out}(\mathit{pk}_{\mathrm{SV}})
      \big\}
    \end{aligned}
  \]
  \textbf{Axioma L\'ogico de Separaci\'on de Custodia:}
  \[
    \begin{aligned}
      \forall t.\; &\neg\,\operatorname{Out}(\mathit{dk}_{\mathrm{SV}})@t \;\land\;
      \Big(\operatorname{SetupSV}(\mathit{ek}, \mathit{pk})@t \implies \\
      &\exists \mathit{dk}, \mathit{sk}.\; \mathit{ek} = \operatorname{ek\_from\_dk}(\mathit{dk}) \land \mathit{pk} = \operatorname{pk}(\mathit{sk})\Big)
    \end{aligned}
  \]
\end{cajaLogica}

\subsection{Regla 3: Registro\_Ciudadano\_Padron}

\begin{lstlisting}[caption={La AR registra al votante y emite su credencial firmada.}]
rule Registro_Ciudadano_Padron:
    let
        pk_V       = pk(~sk_V)
        credencial = <~idV, pk_V, ~ts_reg>   // Paquete de identidad
        firma_AR   = sign(credencial, ~sk_AR) // Firmada (ML-DSA-65)
    in
    [ Fr(~sk_V), Fr(~ts_reg), Fr(~idV)
    , !ClavePrivadaAR('AR', ~sk_AR)
    ]
  --[ RegistroVotante(~idV, pk_V) ]->
    [ !CredencialVotante(~idV, ~sk_V, credencial, firma_AR)
    , !PadronElectoral(~idV, pk_V)
    ]
\end{lstlisting}

\textbf{Nota:} El \textit{timestamp} \texttt{\textasciitilde ts\_reg}
es fresco y \'unico, lo que impide reutilizar credenciales de otras sesiones.

\begin{cajaLogica}{Sem\'antica MSR y L\'ogica Lineal de Registro\_Ciudadano\_Padron}
  \textbf{Transici\'on de Estado MSR (con Hecho Persistente de Autoridad):}
  \[
    \begin{aligned}
      \mathcal{S} \cup \{ \operatorname{Fr}(\mathit{sk}_V), \operatorname{Fr}(\mathit{ts}_{\mathrm{reg}}), \operatorname{Fr}(\mathit{idV}) \}
      &\xrightarrow{\;\operatorname{RegistroVotante}(\mathit{idV}, \operatorname{pk}(\mathit{sk}_V))\;} \\
      \mathcal{S} \cup \big\{
        &!\operatorname{CredencialVotante}(\mathit{idV}, \mathit{sk}_V, \mathit{cred}, \sigma_{\mathrm{AR}}),\,
        !\operatorname{PadronElectoral}(\mathit{idV}, \operatorname{pk}(\mathit{sk}_V))
      \big\}
    \end{aligned}
  \]
  \par\smallskip\noindent\textit{Condicionado a la premisa persistente:} $!\operatorname{ClavePrivadaAR}(\text{'AR'}, \mathit{sk}_{\mathrm{AR}}) \in \mathcal{S}$.\par\medskip
  \textbf{Predicado de Vinculaci\'on de Credencial y Firma ML-DSA-65:}
  \[
    \mathit{cred} = \langle \mathit{idV}, \operatorname{pk}(\mathit{sk}_V), \mathit{ts}_{\mathrm{reg}} \rangle \;\land\;
    \sigma_{\mathrm{AR}} = \operatorname{sign}(\mathit{cred}, \mathit{sk}_{\mathrm{AR}})
  \]
\end{cajaLogica}

\subsection{Regla 4: Emitir\_Voto\_Mesa}

\begin{lstlisting}[caption={El votante cifra su voto con ML-KEM y lo firma con ML-DSA.}]
rule Emitir_Voto_Mesa:
    let
        ct_kem       = encaps(ek_SV, ~k_efimera)   // ML-KEM-768
        voto_cifrado = senc(~m_voto, ~k_efimera)   // AES-256-GCM
        boleta       = <ct_kem, voto_cifrado, credencial, ~ts_voto>
        firma_V      = sign(boleta, ~sk_V)          // ML-DSA-65
    in
    [ Fr(~k_efimera), Fr(~m_voto), Fr(~ts_voto)
    , !CredencialVotante(~idV, ~sk_V, credencial, firma_AR)
    , !ClavePublicaKEM_SV('SV', ek_SV)
    ]
  --[ VotoEmitido(~idV, ~m_voto)
    , VotoEnviado(~idV, boleta, firma_V)
    , SecretoVoto(~m_voto)   // Marca el voto para el lema de secreto
    ]->
    [ Out(<boleta, firma_V, firma_AR>) ]  // Sale a la red
\end{lstlisting}

\textbf{Acci\'on \texttt{SecretoVoto}:} esta etiqueta marca el
contenido del voto para que el lema de confidencialidad pueda
monitorearlo. Sin esta acci\'on, el lema no sabr\'ia a qu\'e dato
aplicar la garant\'ia de secreto.

\begin{cajaLogica}{Sem\'antica MSR y L\'ogica Lineal de Emitir\_Voto\_Mesa}
  \textbf{Transici\'on de Estado MSR (Cifrado H\'ibrido PQC):}
  \[
    \mathcal{S} \cup \{ \operatorname{Fr}(\mathit{k}_{\text{ef\'im}}), \operatorname{Fr}(\mathit{m}_{\text{voto}}), \operatorname{Fr}(\mathit{ts}_{\text{voto}}) \}
    \xrightarrow{\;\mathbf{A}_{\text{emisi\'on}}\;}
    \mathcal{S} \cup \{ \operatorname{Out}(\langle \mathit{boleta}, \sigma_V, \sigma_{\mathrm{AR}} \rangle) \}
  \]
  donde el vector de acciones temporales registradas en el instante $t$ es:
  \[
    \mathbf{A}_{\text{emisi\'on}} = \big[
      \operatorname{VotoEmitido}(\mathit{idV}, \mathit{m}_{\text{voto}}),\,
      \operatorname{VotoEnviado}(\mathit{idV}, \mathit{boleta}, \sigma_V),\,
      \operatorname{SecretoVoto}(\mathit{m}_{\text{voto}})
    \big]
  \]
  \textbf{Proposici\'on de Construcci\'on Criptogr\'afica Cifra-Firma:}
  \[
    \begin{aligned}
      \mathit{boleta} &= \big\langle \operatorname{encaps}(\mathit{ek}_{\mathrm{SV}}, \mathit{k}_{\text{ef\'im}}),\,
      \operatorname{senc}(\mathit{m}_{\text{voto}}, \mathit{k}_{\text{ef\'im}}),\, \mathit{credencial},\, \mathit{ts}_{\text{voto}} \big\rangle \\
      \sigma_V &= \operatorname{sign}(\mathit{boleta}, \mathit{sk}_V)
    \end{aligned}
  \]
\end{cajaLogica}

\subsection{Regla 5: Recepcion\_Y\_Deposito\_Anfora}

\begin{lstlisting}[caption={El SV verifica la boleta y la deposita en el anfora.}]
rule Recepcion_Y_Deposito_Anfora:
    let
        credencial  = <idV, pk_V, ts_reg>
        boleta      = <ct_kem, voto_cifrado, credencial, ts_voto>
        recibo_hash = h(<boleta, ~ts_recibo>)
    in
    [ In(<boleta, firma_V, firma_AR>)
    , !ClavePublicaAR('AR', pk_AR)
    , !PadronElectoral(idV, pk_V)
    , Fr(~ts_recibo)
    ]
  --[ Eq(verify(firma_AR, credencial, pk_AR), true) // Firma AR
    , Eq(verify(firma_V, boleta, pk_V), true)       // Firma Votante
    , VotoUnico(idV)                                // Anti-duplicado
    , VotoAceptado(idV, boleta, firma_V)
    , VotoRegistrado(idV, ct_kem, voto_cifrado)
    ]->
    [ AnforaDigital(idV, ct_kem, voto_cifrado,
                   pk_V, credencial, firma_V, ts_voto)
    , Out(recibo_hash)
    ]
\end{lstlisting}

\textbf{Doble verificaci\'on:} La acci\'on \texttt{Eq(verify(...), true)}
combinada con la restricci\'on \texttt{Verificacion\_Exitosa} garantiza
que Tamarin solo considere trazas donde ambas firmas son v\'alidas.

\begin{cajaLogica}{Sem\'antica MSR y L\'ogica Lineal de Recepci\'on y Dep\'osito en \'Anfora}
  \textbf{Transici\'on de Estado MSR (Consumo de Red y Generaci\'on de Recurso Lineal):}
  \[
    \begin{aligned}
      \mathcal{S} \cup \{ \operatorname{In}(\langle \mathit{boleta}, \sigma_V, \sigma_{\mathrm{AR}} \rangle),\, \operatorname{Fr}(\mathit{ts}_{\text{recibo}}) \}
      &\xrightarrow{\;\mathbf{A}_{\text{rec}}\;} \\
      \mathcal{S} \cup \big\{ \operatorname{AnforaDigital}(\dots),\,
      &\operatorname{Out}(\operatorname{h}(\langle \mathit{boleta}, \mathit{ts}_{\text{recibo}} \rangle)) \big\}
    \end{aligned}
  \]
  \textbf{Guardas L\'ogicas de Verificaci\'on Ecuacional:}
  \[
    \operatorname{verify}(\sigma_{\mathrm{AR}}, \mathit{credencial}, \mathit{pk}_{\mathrm{AR}}) =_E \mathrm{true}
    \quad \land \quad
    \operatorname{verify}(\sigma_V, \mathit{boleta}, \mathit{pk}_V) =_E \mathrm{true}
  \]
  \textbf{Naturaleza Lineal del Recurso \'Anfora:}
  \[
    \operatorname{AnforaDigital} \in \mathcal{F}_{\mathrm{lineal}} \implies
    \text{el hecho es consumido al abrirse el \'anfora, impidiendo doble conteo.}
  \]
\end{cajaLogica}

\subsection{Regla 6: Apertura\_Anfora\_Y\_Escrutinio}

\begin{lstlisting}[caption={La AC desencapsula el secreto KEM y descifra el voto.}]
rule Apertura_Anfora_Y_Escrutinio:
    let
        k_recuperada = decaps(~dk_SV, ct_kem)
        voto_plano   = sdec(voto_cifrado, k_recuperada)
    in
    [ AnforaDigital(idV, ct_kem, voto_cifrado,
                   pk_V, credencial, firma_V, ts_voto)
    , !ClaveDesencapsulamientoAC('AC', ~dk_SV)
    ]
  --[ VotoContado(idV, voto_plano) ]->
    [ Out(<idV, voto_plano>) ]
\end{lstlisting}

\begin{cajaLogica}{Sem\'antica MSR y L\'ogica Lineal de Apertura y Escrutinio}
  \textbf{Transici\'on de Estado MSR (Consumo del Recurso Lineal del \'Anfora):}
  \[
    \mathcal{S} \cup \{ \operatorname{AnforaDigital}(\mathit{idV}, \mathit{ct}_{\text{kem}}, \mathit{voto\_cifrado}, \dots) \}
    \xrightarrow{\;\operatorname{VotoContado}(\mathit{idV}, \mathit{voto\_plano})\;}
    \mathcal{S} \cup \{ \operatorname{Out}(\langle \mathit{idV}, \mathit{voto\_plano} \rangle) \}
  \]
  \par\smallskip\noindent\textit{Condicionado a la premisa persistente:} $!\operatorname{ClaveDesencapsulamientoAC}(\text{'AC'}, \mathit{dk}_{\mathrm{SV}}) \in \mathcal{S}$.\par\medskip
  \textbf{Ecuaciones de Reducci\'on Simb\'olica del Escrutinio:}
  \[
    \mathit{k} = \operatorname{decaps}(\mathit{dk}_{\mathrm{SV}}, \operatorname{encaps}(\operatorname{ek\_from\_dk}(\mathit{dk}_{\mathrm{SV}}), \mathit{k}_{\text{ef\'im}})) =_E \mathit{k}_{\text{ef\'im}}
  \]
  \[
    \mathit{voto\_plano} = \operatorname{sdec}(\operatorname{senc}(\mathit{m}_{\text{voto}}, \mathit{k}_{\text{ef\'im}}), \mathit{k}_{\text{ef\'im}}) =_E \mathit{m}_{\text{voto}}
  \]
  \textbf{Proposici\'on de Correcci\'on del Conteo:}
  \[
    \forall \mathit{idV}, \mathit{m}, t.\; \operatorname{VotoContado}(\mathit{idV}, \mathit{m})@t \implies
    \mathit{m} = \mathit{m}_{\text{voto\_original}}
  \]
\end{cajaLogica}

\begin{cajaDef}{Por qu\'e solo la AC puede descifrar}
  La clave \texttt{\textasciitilde dk\_SV} es persistente con el
  agente \texttt{'AC'}: solo la Autoridad de Conteo tiene acceso a ella.
  Sin esta clave, el ciphertext \texttt{ct\_kem} es computacionalmente
  indistinguible de valores aleatorios (seguridad IND-CCA2 de ML-KEM-768).
\end{cajaDef}

\begin{cajaResumen}
  Las seis reglas modelan el ciclo completo de una votaci\'on: configuraci\'on
  de autoridades, registro, emisi\'on, recepci\'on y escrutinio. Los hechos
  persistentes preservan informaci\'on reutilizable; los lineales representan
  recursos que se consumen exactamente una vez.
\end{cajaResumen}

% ============================================================
\chapter{Lemas y Restricciones: Las Propiedades Formales}
% ============================================================

\epigraph{%
  \textit{``A lemma is a precise mathematical statement about the
  behavior of the protocol under adversarial conditions.''}
}{--- Basin, Dreier \& Sasse, 2018}

\section{Restricciones (Axiomas del sistema)}

Las \textbf{restricciones} permiten filtrar el espacio de trazas,
conservando \'unicamente aquellas coherentes con las invariantes del
protocolo.

\begin{cajaSintaxis}{Restricci\'on}
  \texttt{restriction NombreRestriccion:}\\
  \texttt{\ \ \ \ "formula\_de\_primer\_orden"}
\end{cajaSintaxis}

\begin{lstlisting}[caption={Restricciones del protocolo de votaci\'on.}]
// Cada votante solo puede depositar UN voto en el anfora
restriction Unicidad_Voto:
    "All id #i #j.
        VotoUnico(id) @i & VotoUnico(id) @j ==> #i = #j"

// La verificacion de firma siempre es exitosa cuando se usa Eq
restriction Verificacion_Exitosa:
    "All x y #i. Eq(x, y) @i ==> x = y"
\end{lstlisting}

\begin{cajaDef}{F\'ormulas temporales en Tamarin}
  Las f\'ormulas de los lemas y restricciones usan la siguiente sintaxis:
  \begin{itemize}[itemsep=4pt]
    \item \texttt{Accion(t) @\#i}: la acci\'on \texttt{Accion(t)}
      ocurri\'o en el instante $\#i$.
    \item \texttt{\#i < \#j}: el instante $\#i$ es anterior a $\#j$.
    \item \texttt{All x \#i. P(x) @\#i ==> Q}: para todo $x$ e
      instante $\#i$ en que ocurre $P(x)$, se cumple $Q$.
    \item \texttt{Ex x \#i. P(x) @\#i}: existe alg\'un $x$ e
      instante $\#i$ en que ocurre $P(x)$.
    \item \texttt{not(...)}: negaci\'on l\'ogica.
    \item \texttt{K(t)}: el atacante conoce el t\'ermino $t$.
  \end{itemize}
\end{cajaDef}

\section{Tipos de lemas}

\begin{table}[H]
\centering
\caption{Tipos de lemas en Tamarin Prover.}
\label{tab:tipos-lemas}
\begin{tabularx}{\textwidth}{>{\bfseries\ttfamily}p{3.5cm} X X}
\toprule
\normalfont Tipo & Verificaci\'on & Cu\'ando usar \\
\midrule
exists-trace & Se verifica si \emph{existe} al menos una traza que
  satisface la f\'ormula. & Sanity checks: probar que el protocolo
  puede completarse. \\[6pt]
all-traces & Se verifica si la f\'ormula se cumple en \emph{todas}
  las trazas posibles. & Propiedades de seguridad: confidencialidad,
  autenticidad, integridad. \\
\bottomrule
\end{tabularx}
\end{table}

\section{Los cuatro lemas del protocolo PQC}

\subsection{Lema 1: Ejecutabilidad (Sanity Check)}

\begin{lstlisting}[caption={Lema de ejecutabilidad del protocolo completo.}]
lemma ejecutabilidad:
    exists-trace
    "Ex idV m blt sig #t1 #t2 #t3.
        VotoEmitido(idV, m)         @#t1
      & VotoAceptado(idV, blt, sig) @#t2
      & VotoContado(idV, m)         @#t3
      & #t1 < #t2
      & #t2 < #t3"
\end{lstlisting}

\begin{cajaLogica}{Proposici\'on L\'ogica de Ejecutabilidad (Existencia de Modelo)}
  \[
    \exists \tau \in \operatorname{Traces}(\mathcal{P}),\; \exists \mathit{idV}, m, \mathit{blt}, \mathit{sig} \in \mathcal{T},\; \exists t_1, t_2, t_3 \in \mathbb{N}:
  \]
  \[
    \begin{aligned}
      \tau, t_1 &\models \operatorname{VotoEmitido}(\mathit{idV}, m) \\
      \land\quad \tau, t_2 &\models \operatorname{VotoAceptado}(\mathit{idV}, \mathit{blt}, \mathit{sig}) \\
      \land\quad \tau, t_3 &\models \operatorname{VotoContado}(\mathit{idV}, m) \\
      \land\quad t_1 &< t_2 < t_3
    \end{aligned}
  \]
  \textbf{Interpretaci\'on en Teor\'ia de Modelos:} Esta f\'ormula es de tipo \emph{alcanzabilidad existencial}. Demuestra formalmente que existe al menos una traza admisible $\tau$ en la que el protocolo se ejecuta de principio a fin, garantizando la ausencia de interbloqueos (\emph{deadlocks}) o precondiciones inalcanzables.
\end{cajaLogica}

\textbf{Prop\'osito:} verificar que el protocolo es \textit{ejecutable},
es decir, que existe al menos una historia de ejecuci\'on en la que
todo funciona correctamente en el orden esperado. Si este lema falla,
el modelo est\'a mal construido.

\subsection{Lema 2: Confidencialidad del Voto}

\begin{lstlisting}[caption={Confidencialidad: el atacante Dolev-Yao no aprende el voto.}]
lemma secreto_voto:
    all-traces
    "All m #t.
        SecretoVoto(m) @#t
        ==>
        not(Ex #t2. K(m) @#t2)"
\end{lstlisting}

\begin{cajaLogica}{Proposici\'on L\'ogica del Secreto del Sufragio (Invariante Universal)}
  \[
    \forall \tau \in \operatorname{Traces}(\mathcal{P}),\; \forall m \in \mathcal{T},\; \forall t \in \mathbb{N}:
  \]
  \[
    \tau, t \models \operatorname{SecretoVoto}(m) \implies \neg\,\Big(\exists t_2 \in \mathbb{N}:\; \tau, t_2 \models \mathrm{K}(m)\Big)
  \]
  \textbf{Equivalente por Contraposici\'on:}
  \[
    \exists t_2 \in \mathbb{N}:\; \tau, t_2 \models \mathrm{K}(m) \implies \neg\,\Big(\exists t \in \mathbb{N}:\; \tau, t \models \operatorname{SecretoVoto}(m)\Big)
  \]
  \textbf{Cadena Deductiva Simb\'olica:}
  \[
    \mathrm{K}(m) \iff \Big(\mathrm{K}(\operatorname{senc}(m, k_{\text{ef\'im}})) \;\land\; \mathrm{K}(k_{\text{ef\'im}})\Big)
  \]
  \[
    \mathrm{K}(k_{\text{ef\'im}}) \iff \Big(\mathrm{K}(\operatorname{encaps}(\mathit{ek}_{\mathrm{SV}}, k_{\text{ef\'im}})) \;\land\; \mathrm{K}(\mathit{dk}_{\mathrm{SV}})\Big)
  \]
  Dado que $\forall t.\; \neg\,\mathrm{K}(\mathit{dk}_{\mathrm{SV}})@t$ (pues $\mathit{dk}_{\mathrm{SV}}$ nunca es publicado en un hecho $\operatorname{Out}$), se concluye:
  \[
    \forall t.\; \neg\,\mathrm{K}(k_{\text{ef\'im}})@t \implies \forall t_2.\; \neg\,\mathrm{K}(m)@t_2 \quad \blacksquare
  \]
\end{cajaLogica}

\textbf{Prop\'osito:} garantizar el \textbf{secreto del sufragio}.
El atacante controla la red y puede interceptar todos los mensajes,
incluyendo la boleta cifrada. Sin embargo, sin la clave \texttt{dk\_SV}
(que nunca circula por la red), el contenido del voto es inasequible.

\subsection{Lema 3: Integridad del Sufragio}

\begin{lstlisting}[caption={Integridad: todo voto contado fue leg\'itimamente emitido.}]
lemma integridad_voto:
    all-traces
    "All idV m #tc.
        VotoContado(idV, m) @#tc
        ==>
        Ex #te. VotoEmitido(idV, m) @#te & #te < #tc"
\end{lstlisting}

\begin{cajaLogica}{Proposici\'on L\'ogica de Integridad y Correspondencia Causal}
  \[
    \forall \tau \in \operatorname{Traces}(\mathcal{P}),\; \forall \mathit{idV}, m \in \mathcal{T},\; \forall t_c \in \mathbb{N}:
  \]
  \[
    \begin{aligned}
      \tau, t_c \models \operatorname{VotoContado}(\mathit{idV}, m) \implies
      \exists t_e \in \mathbb{N}:\; \Big( &\tau, t_e \models \operatorname{VotoEmitido}(\mathit{idV}, m) \\
      &\land\; t_e < t_c \Big)
    \end{aligned}
  \]
  \textbf{Significado Sem\'antico:} Ning\'un voto puede emerger en el escrutinio de la AC en el paso $t_c$ a menos que exista un evento causal leg\'itimo previo de emisi\'on por el votante $\mathit{idV}$ en $t_e < t_c$ con el mismo contenido de sufragio $m$. Descarta formalmente la inyecci\'on de votos espurios o manipulaciones en tr\'ansito.
\end{cajaLogica}

\textbf{Prop\'osito:} garantizar la \textbf{correspondencia} entre
voto emitido y voto contado. Ning\'un voto puede aparecer en el
escrutinio sin haber sido emitido leg\'itimamente. Esto descarta votos
fantasma y alteraciones en tr\'ansito.

\subsection{Lema 4: Autenticidad e Injective Agreement}

\begin{lstlisting}[caption={Autenticidad: injective agreement entre env\'io y aceptaci\'on.}]
lemma autenticidad_votante:
    all-traces
    "All idV blt sig #ta.
        VotoAceptado(idV, blt, sig) @#ta
        ==>
        ( Ex #ts.
            VotoEnviado(idV, blt, sig) @#ts
          & #ts < #ta
          & not( Ex #ta2.
                VotoAceptado(idV, blt, sig) @#ta2
              & not(#ta = #ta2) ) )"
\end{lstlisting}

\begin{cajaLogica}{Proposici\'on L\'ogica de Injective Agreement (Autenticidad y Anti-Replay)}
  \[
    \forall \tau \in \operatorname{Traces}(\mathcal{P}),\; \forall \mathit{idV}, \mathit{blt}, \mathit{sig} \in \mathcal{T},\; \forall t_a \in \mathbb{N}:
  \]
  \[
    \begin{aligned}
      \tau, t_a \models \operatorname{VotoAceptado}(\mathit{idV}, \mathit{blt}, \mathit{sig}) \implies
      &\Big( \exists t_s < t_a:\; \tau, t_s \models \operatorname{VotoEnviado}(\mathit{idV}, \mathit{blt}, \mathit{sig}) \Big) \\
      &\land\; \neg\,\Big( \exists t_a' \neq t_a:\; \tau, t_a' \models \operatorname{VotoAceptado}(\mathit{idV}, \mathit{blt}, \mathit{sig}) \Big)
    \end{aligned}
  \]
  \textbf{Descomposici\'on Te\'orica del Acuerdo Inyectivo de Lowe:}
  \begin{itemize}[itemsep=2pt]
    \item \textbf{Autenticidad (Non-injective Agreement):} Todo voto aceptado fue efectivamente enviado por el votante aut\'entico $\mathit{idV}$.
    \item \textbf{Inyectividad Estricta (Anti-replay):} La relaci\'on binaria entre emisiones y aceptaciones es inyectiva; no existe ninguna traza donde una misma boleta sea aceptada m\'as de una vez en instantes temporales distintos ($t_a' \neq t_a$).
  \end{itemize}
\end{cajaLogica}

\textbf{Prop\'osito:} el \textit{injective agreement} garantiza
\textbf{dos propiedades simult\'aneamente}: (1)~autenticidad --- la
boleta proviene del votante leg\'itimo; (2)~anti-replay --- la misma
boleta no puede ser aceptada m\'as de una vez.

\begin{cajaResumen}
  Los lemas son el coraz\'on de la verificaci\'on formal. El lema
  \texttt{exists-trace} verifica ejecutabilidad; los de
  \texttt{all-traces} verifican propiedades de seguridad universales.
  Las cuatro propiedades del protocolo PQC cubren los requisitos
  m\'inimos de un sistema electoral: confidencialidad, integridad,
  autenticidad y unicidad del voto.
\end{cajaResumen}

% ────────────────────────────────────────────────────────────
%   PARTE III — EL PROTOCOLO PQC Y SUS RESULTADOS
% ────────────────────────────────────────────────────────────
\part{Protocolo PQC: An\'alisis y Resultados}

% ============================================================
\chapter{El Protocolo de Votaci\'on Electr\'onica Post-Cu\'antica}
% ============================================================

\epigraph{%
  \textit{``A quantum computer powerful enough to break today's
  cryptography may be a decade away. But the data encrypted today
  may still need to be protected in a decade.''}
}{--- NIST, Special Publication 800-208}

\section{Contexto: criptograf\'ia post-cu\'antica}

Los algoritmos criptogr\'aficos tradicionales (RSA, ECDSA, Diffie-Hellman)
basan su seguridad en la dificultad de ciertos problemas matem\'aticos.
El algoritmo de Shor, ejecutable en una computadora cu\'antica,
puede resolver estos problemas en tiempo polinomial, rompiendo todos
estos esquemas.

\medskip

El NIST respondi\'o con un proceso de estandarizaci\'on de algoritmos
post-cu\'anticos, concluido en 2024 con:
\begin{itemize}[itemsep=3pt]
  \item \textbf{FIPS~203 (ML-KEM)}: mecanismo de encapsulamiento de
    claves basado en ret\'iculos modulares.
  \item \textbf{FIPS~204 (ML-DSA)}: esquema de firma digital basado
    en ret\'iculos modulares.
\end{itemize}

\section{Arquitectura del protocolo}

\begin{figure}[H]
\centering
\begin{tikzpicture}[
  entity/.style={draw=TamarinBlue, fill=TamarinLight!70,
                 rounded corners=8pt, inner sep=10pt,
                 text width=3cm, align=center, font=\small\bfseries},
  arrow/.style={-Stealth, thick, TamarinBlue},
  lbl/.style={font=\scriptsize\itshape, text=gray, align=center}
]
  \node[entity] (AR) at (0,4)   {AR\\{\normalfont\tiny Aut. de Registro}};
  \node[entity] (V)  at (0,0)   {V\\{\normalfont\tiny Votante}};
  \node[entity] (SV) at (8,2)   {SV\\{\normalfont\tiny Servidor Votaci\'on}};
  \node[entity] (AC) at (4,-2.5){AC\\{\normalfont\tiny Aut. de Conteo}};

  \draw[arrow] (AR) -- node[left, lbl]{(b) Credencial firmada} (V);
  \draw[arrow] (V)  -- node[above, lbl]{(c) Boleta cifrada + firmada} (SV);
  \draw[arrow, dashed] (SV) -- node[right, lbl]{dk\_SV\\(fuera de banda)} (AC);
  \draw[arrow] (SV) -- node[above left, lbl]{\'Anfora Digital\\(tras cierre)} (AC);
  \draw[arrow, dashed, gray] (AR) -- node[above, lbl]{(a) Config.} (SV);
\end{tikzpicture}
\caption{Arquitectura y flujo de mensajes del protocolo de votaci\'on PQC.}
\label{fig:arquitectura}
\end{figure}

\section{Las entidades del protocolo}

\begin{table}[H]
\centering
\caption{Roles del protocolo de votaci\'on electr\'onica PQC.}
\label{tab:roles}
\begin{tabular}{>{\bfseries\color{TamarinBlue}}p{1.2cm}
                >{\bfseries}p{4cm}
                p{7.5cm}}
\toprule
Rol & Entidad & Responsabilidad \\
\midrule
AR & Autoridad de Registro &
  Verifica identidades, genera claves del votante y emite credenciales
  firmadas con ML-DSA-65. \\[4pt]
SV & Servidor de Votaci\'on &
  Recepciona boletas, verifica firmas, previene doble voto y custodia
  el \'Anfora Digital. \\[4pt]
AC & Autoridad de Conteo &
  Custodia exclusivamente \texttt{dk\_SV}. Al cierre electoral, abre
  el \'anfora y totaliza los votos. \\[4pt]
V  & Votante / Ciudadano &
  Cifra su voto con ML-KEM-768, firma la boleta con ML-DSA-65 y la
  env\'ia al SV. \\
\bottomrule
\end{tabular}
\end{table}

\section{Las seis fases del protocolo}

\begin{enumerate}[label=\textbf{Fase \Alph*:}, leftmargin=*, itemsep=8pt]

  \item \textbf{Configuraci\'on de la AR.} La Autoridad de Registro
    genera su par de claves ML-DSA-65 (\texttt{sk\_AR}, \texttt{pk\_AR})
    y publica \texttt{pk\_AR}.

  \item \textbf{Configuraci\'on del SV.} El Servidor de Votaci\'on genera
    su par ML-KEM-768 (\texttt{dk\_SV}, \texttt{ek\_SV}) y su par
    ML-DSA-65 (\texttt{sk\_SV}, \texttt{pk\_SV}). Entrega
    \texttt{dk\_SV} exclusivamente a la AC fuera de banda y publica
    las claves p\'ublicas.

  \item \textbf{Registro en el Padr\'on.} La AR verifica la identidad
    del votante, genera su par de claves, construye la credencial
    $\langle idV, pk_V, ts_{reg} \rangle$, la firma y la entrega
    junto con \texttt{sk\_V}.

  \item \textbf{Emisi\'on del Voto.} El votante: (1)~genera
    \texttt{\textasciitilde k\_efimera}; (2)~encapsula:
    $ct_{kem} = \texttt{encaps}(ek_{SV}, k_{ef})$; (3)~cifra:
    $vc = \texttt{senc}(m_{voto}, k_{ef})$; (4)~construye boleta;
    (5)~firma; (6)~env\'ia al SV.

  \item \textbf{Verificaci\'on y Dep\'osito (SV).} El SV verifica las
    dos firmas, comprueba unicidad del votante, deposita en el \'Anfora
    Digital y emite recibo $h(\langle blt, ts_{recibo}\rangle)$.

  \item \textbf{Escrutinio (AC).} Tras el cierre electoral, la AC
    recupera $k = \texttt{decaps}(dk_{SV}, ct_{kem})$, descifra
    $m = \texttt{sdec}(vc, k)$ y totaliza los resultados.

\end{enumerate}

\section{Propiedades de seguridad requeridas}

\begin{table}[H]
\centering
\caption{Propiedades de seguridad del protocolo y su fundamento.}
\label{tab:propiedades}
\begin{tabularx}{\textwidth}{>{\bfseries}p{3cm} X p{3.5cm}}
\toprule
Propiedad & Descripci\'on & Mecanismo \\
\midrule
Confidencialidad &
  El contenido del voto no es accesible salvo a la AC, y solo
  tras el cierre electoral. &
  ML-KEM-768 + AES-256-GCM \\[5pt]
Autenticidad &
  Solo votantes registrados pueden emitir boletas v\'alidas. &
  ML-DSA-65 \\[5pt]
Integridad &
  El voto contado es exactamente el voto emitido. &
  Firma boleta + ecuaci\'on KEM \\[5pt]
Unicidad &
  Cada votante puede emitir exactamente un voto. &
  Restricci\'on \texttt{VotoUnico} \\[5pt]
Anti-replay &
  Una boleta v\'alida no puede ser reutilizada. &
  Timestamp fresco + Injective Agr. \\
\bottomrule
\end{tabularx}
\end{table}

% ============================================================
\chapter{Ejecuci\'on y Resultados de la Verificaci\'on}
% ============================================================

\epigraph{%
  \textit{``verified doesn't mean perfect. It means: within the model,
  under these assumptions, no attack exists. That is both very strong
  and very precise.''}
}{--- Simon Meier, ETH Z\"urich}

\section{Requisitos del sistema}

\begin{table}[H]
\centering
\caption{Requisitos del sistema para ejecutar Tamarin.}
\label{tab:requisitos}
\begin{tabular}{>{\bfseries}p{4cm} p{8cm}}
\toprule
Componente & Versi\'on / Notas \\
\midrule
Tamarin Prover & \texttt{tamarin-prover-1.8.x} o superior \\[3pt]
Maude & 3.5.1 (incluido con el instalador de Tamarin) \\[3pt]
Sistema operativo & Linux (recomendado), macOS, Windows (WSL2) \\[3pt]
RAM & M\'inimo 4~GB; recomendado 8~GB o m\'as \\[3pt]
Tiempo de verificaci\'on & $\approx 30$~s (protocolo PQC, 4~lemas) \\
\bottomrule
\end{tabular}
\end{table}

\section{Comandos de ejecuci\'on}

\begin{lstlisting}[style=BashStyle,
                   caption={Comandos para verificar el protocolo PQC.}]
# Verificar todos los lemas automaticamente
tamarin-prover --prove tamarin/protocolo.spthy

# Verificar un solo lema
tamarin-prover --prove=secreto_voto tamarin/protocolo.spthy

# Modo interactivo con interfaz web (GUI)
# Abrir navegador en http://localhost:3001
tamarin-prover interactive tamarin/protocolo.spthy

# Verificacion con multiples nucleos (4 hilos)
tamarin-prover --prove +RTS -N4 -RTS tamarin/protocolo.spthy
\end{lstlisting}

\section{Salida del motor Maude}

Tras iniciar la verificaci\'on, Tamarin reporta el progreso:

\begin{lstlisting}[style=BashStyle,
                   caption={Salida del motor Maude (reporte\_tamarin.txt).}]
maude tool: 'maude'
 checking version: 3.5.1. OK.
 checking installation: OK.
[Theory ProtocoloVotacionPQC] Theory loaded
[Theory ProtocoloVotacionPQC] Theory translated
[Theory ProtocoloVotacionPQC] Derivation checks started
[Theory ProtocoloVotacionPQC] Derivation checks ended
[Theory ProtocoloVotacionPQC] Theory closed
[Saturating Sources] Step 1 (Max 5)
[Saturating Sources] Done
\end{lstlisting}

\section{Resultados de la verificaci\'on formal}

\begin{cajaResultado}{\Large Resultados: 4 / 4 Lemas Verificados}
\begin{table}[H]
\centering
\begin{tabular}{>{\ttfamily}p{5cm}
                p{5.5cm}
                >{\bfseries\color{AccentGreen}}p{2.5cm}}
\toprule
\normalfont Lema &
\normalfont Propiedad garantizada &
\normalfont Estado \\
\midrule
ejecutabilidad &
  Traza completa ejecutable (sanity check) &
  VERIFIED \\[5pt]
secreto\_voto &
  Confidencialidad del sufragio &
  VERIFIED \\[5pt]
integridad\_voto &
  Correspondencia voto emitido/contado &
  VERIFIED \\[5pt]
autenticidad\_votante &
  Anti-replay (Injective Agreement) &
  VERIFIED \\
\bottomrule
\end{tabular}
\end{table}
\end{cajaResultado}

\bigskip

\begin{cajaDef}{Qu\'e significa \texttt{VERIFIED}}
  Cuando Tamarin devuelve \texttt{verified} para un lema de tipo
  \texttt{all-traces}, ha explorado \textbf{el espacio completo de
  trazas} del protocolo (posiblemente infinito, gracias a las t\'ecnicas
  de saturaci\'on de fuentes e inducci\'on sobre el n\'umero de sesiones)
  y no ha encontrado ninguna traza que viole la propiedad declarada.

  Para el lema \texttt{exists-trace}, ha encontrado al menos una traza
  concreta que satisface la f\'ormula.

  En ambos casos, el resultado es una \textbf{demostraci\'on matem\'atica}
  dentro del modelo declarado, no una prueba emp\'irica.
\end{cajaDef}

\section{Limitaciones del modelo}

\begin{table}[H]
\centering
\caption{Alcance y limitaciones de la verificaci\'on formal.}
\label{tab:limites}
\begin{tabularx}{\textwidth}{X X}
\toprule
\textbf{Lo que S\'I garantiza} & \textbf{Lo que NO cubre} \\
\midrule
Seguridad en el modelo simb\'olico Dolev-Yao &
  Seguridad computacional de ML-KEM / ML-DSA \\[4pt]
Resistencia a ataques de red (replay, MITM) &
  Ataques de canal lateral (timing, power) \\[4pt]
Imposibilidad de doble voto en el modelo &
  Errores en la implementaci\'on Python/C \\[4pt]
Autenticidad e integridad simb\'olica &
  Ataques al sistema operativo o hardware \\[4pt]
Confidencialidad del voto ante Dolev-Yao &
  Coerci\'on del votante o compra de voto \\
\bottomrule
\end{tabularx}
\end{table}

\begin{cajaAviso}{La verificaci\'on no reemplaza la implementaci\'on segura}
  La verificaci\'on formal es un \textbf{complemento} a las pruebas
  unitarias, las auditor\'ias de c\'odigo y el an\'alisis de seguridad
  computacional. Garantiza que el \textit{dise\~no} del protocolo es
  correcto; no garantiza que la \textit{implementaci\'on} lo sea.
\end{cajaAviso}

\section{Interpretaci\'on de los resultados en el contexto de la tesis}

\begin{enumerate}[leftmargin=*, itemsep=8pt]
  \item \textbf{Rigor cient\'ifico:} la demostraci\'on formal eleva el
    trabajo m\'as all\'a de una implementaci\'on experimental, situ\'andolo
    en el nivel de la verificaci\'on matem\'atica.
  \item \textbf{Confianza en el dise\~no:} el modelo Dolev-Yao es el
    est\'andar de la industria para el an\'alisis de protocolos. Haber
    verificado el protocolo bajo este modelo garantiza que no existen
    vulnerabilidades estructurales conocidas en el dise\~no.
  \item \textbf{Alineaci\'on con est\'andares:} las primitivas ML-KEM-768
    y ML-DSA-65 est\'an estandarizadas por el NIST (FIPS~203 y FIPS~204),
    y su abstracci\'on simb\'olica en Tamarin es coherente con sus
    propiedades formalmente definidas.
\end{enumerate}

\begin{cajaResumen}
  Los 4/4 lemas verificados demuestran formalmente que el protocolo de
  votaci\'on PQC es confidencial, \'integro, aut\'entico y resistente a
  ataques de repetici\'on bajo el modelo Dolev-Yao. La verificaci\'on fue
  realizada con Tamarin Prover sobre el motor Maude 3.5.1.
\end{cajaResumen}

% ── APÉNDICES ───────────────────────────────────────────────
\appendix

% ============================================================
\chapter{C\'odigo Fuente Completo: \texttt{protocolo.spthy}}
% ============================================================

A continuaci\'on se reproduce el archivo completo de especificaci\'on
formal para referencia:

\begin{lstlisting}[caption={Archivo completo protocolo.spthy},
                   label={lst:full-protocol}]
theory ProtocoloVotacionPQC
begin

/*
================================================================================
  TESIS DE GRADO: PROTOCOLO DE VOTACION ELECTRONICA BASADO EN CRIPTOGRAFIA
                  POST-CUANTICA DE RETICULOS (ML-KEM-768 + ML-DSA-65)
  Universidad Mayor de San Andres (UMSA) -- 2026
================================================================================
*/

builtins: dest-signing, symmetric-encryption

functions:
    encaps/2,
    decaps/2,
    ek_from_dk/1,
    h/1

equations:
    decaps(dk, encaps(ek_from_dk(dk), r)) = r

restriction Unicidad_Voto:
    "All id #i #j. VotoUnico(id) @i & VotoUnico(id) @j ==> #i = #j"

restriction Verificacion_Exitosa:
    "All x y #i. Eq(x, y) @i ==> x = y"

rule Setup_Autoridad_Registro:
    let  pk_AR = pk(~sk_AR)  in
    [ Fr(~sk_AR) ]
  --[ SetupAR(pk_AR) ]->
    [ !ClavePrivadaAR('AR', ~sk_AR)
    , !ClavePublicaAR('AR', pk_AR)
    , Out(pk_AR)
    ]

rule Setup_Servidor_Votacion:
    let
        ek_SV = ek_from_dk(~dk_SV)
        pk_SV = pk(~sk_SV)
    in
    [ Fr(~dk_SV), Fr(~sk_SV) ]
  --[ SetupSV(ek_SV, pk_SV) ]->
    [ !ClaveDesencapsulamientoAC('AC', ~dk_SV)
    , !ClavePublicaKEM_SV('SV', ek_SV)
    , !ClavePublicaFirma_SV('SV', pk_SV)
    , Out(ek_SV)
    , Out(pk_SV)
    ]

rule Registro_Ciudadano_Padron:
    let
        pk_V       = pk(~sk_V)
        credencial = <~idV, pk_V, ~ts_reg>
        firma_AR   = sign(credencial, ~sk_AR)
    in
    [ Fr(~sk_V), Fr(~ts_reg), Fr(~idV)
    , !ClavePrivadaAR('AR', ~sk_AR)
    ]
  --[ RegistroVotante(~idV, pk_V) ]->
    [ !CredencialVotante(~idV, ~sk_V, credencial, firma_AR)
    , !PadronElectoral(~idV, pk_V)
    ]

rule Emitir_Voto_Mesa:
    let
        ct_kem       = encaps(ek_SV, ~k_efimera)
        voto_cifrado = senc(~m_voto, ~k_efimera)
        boleta       = <ct_kem, voto_cifrado, credencial, ~ts_voto>
        firma_V      = sign(boleta, ~sk_V)
    in
    [ Fr(~k_efimera), Fr(~m_voto), Fr(~ts_voto)
    , !CredencialVotante(~idV, ~sk_V, credencial, firma_AR)
    , !ClavePublicaKEM_SV('SV', ek_SV)
    ]
  --[ VotoEmitido(~idV, ~m_voto)
    , VotoEnviado(~idV, boleta, firma_V)
    , SecretoVoto(~m_voto)
    ]->
    [ Out(<boleta, firma_V, firma_AR>) ]

rule Recepcion_Y_Deposito_Anfora:
    let
        credencial  = <idV, pk_V, ts_reg>
        boleta      = <ct_kem, voto_cifrado, credencial, ts_voto>
        recibo_hash = h(<boleta, ~ts_recibo>)
    in
    [ In(<boleta, firma_V, firma_AR>)
    , !ClavePublicaAR('AR', pk_AR)
    , !PadronElectoral(idV, pk_V)
    , Fr(~ts_recibo)
    ]
  --[ Eq(verify(firma_AR, credencial, pk_AR), true)
    , Eq(verify(firma_V, boleta, pk_V), true)
    , VotoUnico(idV)
    , VotoAceptado(idV, boleta, firma_V)
    , VotoRegistrado(idV, ct_kem, voto_cifrado)
    ]->
    [ AnforaDigital(idV, ct_kem, voto_cifrado,
                   pk_V, credencial, firma_V, ts_voto)
    , Out(recibo_hash)
    ]

rule Apertura_Anfora_Y_Escrutinio:
    let
        k_recuperada = decaps(~dk_SV, ct_kem)
        voto_plano   = sdec(voto_cifrado, k_recuperada)
    in
    [ AnforaDigital(idV, ct_kem, voto_cifrado,
                   pk_V, credencial, firma_V, ts_voto)
    , !ClaveDesencapsulamientoAC('AC', ~dk_SV)
    ]
  --[ VotoContado(idV, voto_plano) ]->
    [ Out(<idV, voto_plano>) ]

lemma ejecutabilidad:
    exists-trace
    "Ex idV m blt sig #t1 #t2 #t3.
        VotoEmitido(idV, m)         @#t1
      & VotoAceptado(idV, blt, sig) @#t2
      & VotoContado(idV, m)         @#t3
      & #t1 < #t2
      & #t2 < #t3"

lemma secreto_voto:
    all-traces
    "All m #t.
        SecretoVoto(m) @#t
        ==>
        not(Ex #t2. K(m) @#t2)"

lemma integridad_voto:
    all-traces
    "All idV m #tc.
        VotoContado(idV, m) @#tc
        ==>
        Ex #te. VotoEmitido(idV, m) @#te & #te < #tc"

lemma autenticidad_votante:
    all-traces
    "All idV blt sig #ta.
        VotoAceptado(idV, blt, sig) @#ta
        ==>
        ( Ex #ts.
            VotoEnviado(idV, blt, sig) @#ts
          & #ts < #ta
          & not( Ex #ta2.
                VotoAceptado(idV, blt, sig) @#ta2
              & not(#ta = #ta2) ) )"

end
\end{lstlisting}

% ============================================================
\chapter{Glosario de T\'erminos T\'ecnicos}
% ============================================================

\begin{description}[leftmargin=3.5cm, style=nextline,
                    font=\bfseries\color{TamarinBlue}]

  \item[all-traces] Tipo de lema en Tamarin. La propiedad debe
    cumplirse en \textit{todas} las trazas posibles. Se usa para
    propiedades de seguridad universales.

  \item[\'Anfora Digital] Contenedor l\'ogico (hecho lineal de Tamarin)
    donde el SV deposita las boletas cifradas para su escrutinio.

  \item[Builtin] M\'odulo predefinido de Tamarin con funciones y
    ecuaciones para operaciones criptogr\'aficas est\'andar.

  \item[Caja negra perfecta] Abstracci\'on del modelo simb\'olico:
    una operaci\'on criptogr\'afica que solo puede invertirse con
    el conocimiento correcto (la clave).

  \item[Dolev-Yao] Modelo de atacante que controla completamente el
    canal de red pero no puede romper las primitivas criptogr\'aficas.

  \item[exists-trace] Tipo de lema en Tamarin. Se verifica si
    \textit{existe} al menos una traza que la satisface. Se usa
    para sanity checks.

  \item[Hecho fresco (\textasciitilde)] Variable generada con
    \texttt{Fr(\textasciitilde x)} en Tamarin. Valor aleatorio
    \'unico no reutilizado.

  \item[Hecho lineal] Hecho del multiconjunto de estado que se
    consume al ser usado como premisa.

  \item[Hecho persistente (!)] Hecho que no se consume al usarse;
    puede reutilizarse m\'ultiples veces.

  \item[Injective Agreement] Propiedad que garantiza que por cada
    aceptaci\'on de un mensaje, existe exactamente un env\'io previo
    y \'unico.

  \item[K(t)] Hecho especial de Tamarin: el atacante conoce $t$.

  \item[KEM] Key Encapsulation Mechanism. Esquema que permite
    encapsular un secreto con la clave p\'ublica del destinatario.

  \item[Lemma / Lema] Declaraci\'on formal de una propiedad que
    Tamarin debe probar o refutar.

  \item[ML-DSA] Module Lattice Digital Signature Algorithm.
    Firma digital post-cu\'antica (FIPS~204).

  \item[ML-KEM] Module Lattice Key Encapsulation Mechanism.
    KEM post-cu\'antico (FIPS~203, antes Kyber).

  \item[Multiconjunto de hechos] Representaci\'on del estado global
    del protocolo en Tamarin.

  \item[Restricci\'on] Axioma del sistema en Tamarin. Filtra las
    trazas no coherentes con las invariantes del protocolo.

  \item[Regla de reescritura] Mecanismo b\'asico de transici\'on en
    Tamarin. Consume hechos de las premisas y produce nuevos hechos.

  \item[Sanity check] Verificaci\'on de que el modelo es coherente
    y el protocolo puede completarse (lema \texttt{exists-trace}).

  \item[.spthy] Extensi\'on de archivo de las teor\'ias Tamarin
    (\textit{security protocol theory}).

  \item[Tamarin Prover] Herramienta de verificaci\'on formal de
    protocolos criptogr\'aficos. Desarrollada en ETH Z\"urich.
    Usa Maude como motor de reescritura.

  \item[Traza de ejecuci\'on] Secuencia de estados y acciones que
    describe una posible historia de ejecuci\'on del protocolo.

  \item[Verificaci\'on formal] Demostraci\'on matem\'atica rigurosa
    de que un sistema cumple ciertas propiedades bajo todas las
    condiciones posibles.

\end{description}

% ── BIBLIOGRAFÍA ────────────────────────────────────────────
\bibliographystyle{plain}
\begin{thebibliography}{99}

\bibitem{tamarin2012}
  Schmidt, B., Meier, S., Cremers, C., \& Basin, D. (2012).
  \textit{The TAMARIN Prover for the Symbolic Analysis of Security Protocols}.
  In \textit{Proceedings of CAV 2012}, LNCS, vol.~7358, pp.~696--701. Springer.

\bibitem{cremers2012}
  Cremers, C., \& Mauw, S. (2012).
  \textit{Operational Semantics and Verification of Security Protocols}.
  Springer.

\bibitem{fips203}
  National Institute of Standards and Technology. (2024).
  \textit{FIPS 203 --- Module-Lattice-Based Key-Encapsulation Mechanism Standard}.
  U.S. Department of Commerce.

\bibitem{fips204}
  National Institute of Standards and Technology. (2024).
  \textit{FIPS 204 --- Module-Lattice-Based Digital Signature Standard}.
  U.S. Department of Commerce.

\bibitem{dolevyao1983}
  Dolev, D., \& Yao, A. C. (1983).
  On the security of public key protocols.
  \textit{IEEE Transactions on Information Theory}, 29(2), 198--208.

\bibitem{basin2015}
  Basin, D., Dreier, J., \& Sasse, R. (2015).
  Automated Symbolic Proofs of Observational Equivalence.
  In \textit{Proceedings of CCS 2015}, pp.~1144--1155. ACM.

\bibitem{kyber2018}
  Bos, J., Ducas, L., Kiltz, E., et al. (2018).
  CRYSTALS-Kyber: A CCA-Secure Module-Lattice-Based KEM.
  In \textit{EuroS\&P 2018}, pp.~353--367. IEEE.

\bibitem{dilithium2018}
  Ducas, L., Kiltz, E., Lepoint, T., et al. (2018).
  CRYSTALS-Dilithium: A Lattice-Based Digital Signature Scheme.
  \textit{IACR Transactions on Cryptographic Hardware and Embedded
  Systems}, 2018(1), 238--268.

\bibitem{tamarin-manual}
  The Tamarin Team. (2023).
  \textit{Tamarin Prover Manual}.
  \url{https://tamarin-prover.github.io/manual/}

\bibitem{pqclean}
  PQClean Project. (2024).
  \textit{Clean, portable implementations of post-quantum cryptography}.
  \url{https://github.com/PQClean/PQClean}

\end{thebibliography}

\printindex

\end{document}
